\section{Structure Learning: Constraint-Based and Score-Based Methods}
\label{sec:structure-learning}

Structure learning is the harder of the two belief-network learning problems
introduced in Section~\ref{sec:pgm-fundamentals}: given only a dataset, find
the DAG $G$ that best represents the dependencies between variables, without
assuming the structure in advance. The search space is enormous — for $n$
variables there are $n(n-1)/2$ possible undirected edges and therefore up to
$2^{n(n-1)/2}$ candidate structures — so an exhaustive search is computationally
infeasible beyond a handful of variables, and the literature instead relies
on two broad families of approximate algorithm: constraint-based and score-based learning.

\subsection{Constraint-Based Learning}

Constraint-based algorithms use statistical conditional-independence tests
(e.g.\ Pearson's $\chi^2$ test for categorical data, or Fisher's $Z$-test for
Gaussian data) to directly determine which arcs belong in the graph, rather
than optimising a global score. Every constraint-based algorithm shares the
same three-phase structure, illustrated by the Inferred Causality algorithm described in Section~\ref{sec:inferred-causality}: 

\begin{itemize}
	\item learn each
	node's Markov blanket to prune the candidate structure early;
	\item use the reduced candidate set to identify which undirected arcs (the "skeleton") are present
	\item orient the remaining undirected arcs into a
	final, partially-directed acyclic graph.
\end{itemize}

The appeal of this family is that it directly tests the conditional-independence relationships that a DAG is
meant to represent, and the Markov-blanket pruning step (Section \ref{sec:markov-chain-blanket}) gives a substantial computational saving over testing every possible subset of variables. Its main weakness is sensitivity
to the statistical independence tests themselves: with a finite, noisy sample (as in our 24-variable panel, where some EIA series only begin in 2003),
individual test errors can propagate into structural errors that are hard to correct downstream.

\subsection{Score-Based Learning}

Score-based algorithms instead assign a numerical score
$\mathrm{Score}(G, D)$ to every candidate structure $G$, quantifying how well
$G$ represents the observed dataset $D$, and search for the structure that
maximises it. Writing the score as the posterior probability of the graph
given the data,
\[
  \mathrm{Score}(G, D) = \Pr[G \mid D] \;\propto\; \Pr[D \mid G]\,\Pr[G],
\]
a common and well-behaved choice is the \textbf{Bayesian Information
Criterion (BIC)},
\[
  \mathrm{BIC}(G, D) = \sum_{v \in V} \log \hat p_{v \mid \mathrm{pa}(v)}
    \bigl(X_v \mid \mathrm{pa}(v)\bigr) - \tfrac{1}{2} k_{v \mid \mathrm{pa}(v)}
    \log n,
\]
where $k_{v\mid\mathrm{pa}(v)}$ is the number of free parameters at node $v$
given its parents. The first term rewards structures that fit the data well;
the second penalises structures with more free parameters, which both
discourages overfitting and offsets BIC's bias towards overly simple
structures as the sample size $n$ grows.

Because an exhaustive search over all candidate DAGs remains infeasible even
with a scoring function in hand, score-based learning is typically paired
with a local search heuristic. The most widely used is \textbf{Hill
Climbing}: starting from an initial structure (empty, full, random, or
seeded with expert knowledge), the algorithm repeatedly tries every
single-edge modification available at the current structure — adding an
edge, removing an edge, or reversing an edge's direction — and moves to
whichever neighbouring structure most improves the score. The procedure
halts when no single-edge change improves the score further. Hill Climbing
is greedy and therefore has no guarantee of finding the globally optimal
structure; it can settle at a local maximum depending on its starting point.
Standard mitigations include a tabu list of recently-visited
structures (to discourage immediately undoing the last move) and random restarts from a different starting structure.

\subsection{Trade-offs and Relevance}

Constraint-based and score-based learning make different trade-offs: the
former directly tests the conditional-independence structure the DAG is
supposed to represent, at the cost of sensitivity to individual test errors;
the latter avoids relying on any single statistical test by instead
optimising a global, sample-size-aware score, at the cost of only finding a
local optimum via a greedy search. For this project, we use
\texttt{pgmpy}'s \texttt{HillClimbSearch} with a BIC score as the primary
score-based structure-learning tool, since it scales comfortably to our panel size, requires no manual tuning of independence-test thresholds, and its score explicitly trades off fit against the number of free parameters — a desirable property given how many of our 24 series are
themselves strongly collinear (Figure~\ref{fig:correlation-heatmap}), which
would otherwise encourage a constraint-based test to retain spurious edges.
We separately implement the constraint-based Inferred Causality algorithm by hand in Section~\ref{sec:inferred-causality}, both as a pedagogical
exercise and as a cross-check on the Hill-Climbing result.
